\section{Introduction}

Many clinical AI proposals start with the model. They ask which model is most accurate, which retrieval method should be used, or which interface can make the answer easier to read. Those questions matter, but they are too narrow. In a clinical setting, an AI system also has to fit into an institution, a regulatory environment, a workflow, a risk posture, and a set of professional norms. A technically impressive model can still fail if it creates liability, adds work for clinicians, produces explanations that cannot be audited, or reaches patients in the wrong way.

This paper develops that argument through a case study from MIT's System Design and Management (SDM) program. SDM combines system architecture, systems engineering, and project management in an integrated core built around complex socio-technical problems \cite{mit_sdm_core}. The case is Sanguine, a proposed physician-facing clinical interpretation layer for Mayo Clinic laboratory results. The concept uses an agentic AI clinical librarian to augment physician-facing test results with relevant research, while keeping diagnosis and final clinical judgment with the physician.

The paper is not a claim that Sanguine should be deployed as described. It is a first synthesis of the SDM work into a research-style article. The value of the case is in the design path: how the team narrowed the system boundary, selected architecture decisions, compared alternatives, modeled uncertainty, chose a development approach, and revised the baseline after risk analysis.

The paper makes three contributions. First, it documents a clinical AI architecture case in which governance and release control are treated as first-order design decisions. Second, it shows how tradespace exploration changes when residual risk is added alongside cost and schedule. Third, it extracts a set of lessons for designing AI systems in domains where trust, validation, and institutional accountability are not optional.
